\documentclass[12pt]{article}

\usepackage[utf8]{inputenc}
\usepackage{amsmath, amssymb, amsthm}
\usepackage{geometry}
\geometry{a4paper, margin=2.5cm}

\title{Funcții Analitice și Olomorfe}
\author{}
\date{}

\theoremstyle{definition}
\newtheorem{definition}{Definiție}

\theoremstyle{plain}
\newtheorem{proposition}{Propoziție}

\begin{document}

\maketitle

\begin{definition}
Fie $\Omega \subset \mathbb{C}$ un domeniu.  
O funcție $f:\Omega \to \mathbb{C}$ se numește \emph{analitică} în $\Omega$ dacă pentru orice punct $z_0 \in \Omega$ există un disc $D(z_0,r)$ și o serie de puteri


\[
f(z) = \sum_{n=0}^{\infty} a_n (z - z_0)^n
\]


care converge către $f(z)$ pentru toate $z \in D(z_0,r)$.
\end{definition}

\begin{proposition}
(\textbf{Echivalența dintre analitic și olomorf})  
Fie $\Omega \subset \mathbb{C}$ un domeniu și $f:\Omega \to \mathbb{C}$ o funcție.  
Atunci următoarele afirmații sunt echivalente:
\begin{enumerate}
    \item[(i)] Funcția $f$ este \emph{analitică} în $\Omega$.
    \item[(ii)] Funcția $f$ este \emph{olomorfă} în $\Omega$, adică derivabilă complex în fiecare punct al lui $\Omega$.
\end{enumerate}
Cu alte cuvinte, într-un domeniu din $\mathbb{C}$, analiticitatea și olomorfia coincid.
\end{proposition}

\end{document}
